\section{AMC 10 10B 2004}
        \begin{enumerate}
        	\item (\textit{AMC 10 10B 2004}) Each row of the Misty Moon Amphitheater has $33$ seats. Rows $12$ through $22$ are reserved for a youth club. How many seats are reserved for this club?


 $\mathrm{(A) \ } 297 \qquad \mathrm{(B) \ } 330\qquad \mathrm{(C) \ } 363\qquad \mathrm{(D) \ } 396\qquad \mathrm{(E) \ } 726$


	\item (\textit{AMC 10 10B 2004}) How many two-digit positive integers have at least one $7$ as a digit?


 $\mathrm{(A) \ } 10 \qquad \mathrm{(B) \ } 18\qquad \mathrm{(C) \ } 19 \qquad \mathrm{(D) \ } 20\qquad \mathrm{(E) \ } 30$


	\item (\textit{AMC 10 10B 2004}) At each basketball practice last week, Jenny made twice as many free throws as she made at the previous practice. At her fifth practice she made $48$ free throws. How many free throws did she make at the first practice?


 $\mathrm{(A) \ } 3 \qquad \mathrm{(B) \ } 6 \qquad \mathrm{(C) \ } 9 \qquad \mathrm{(D) \ } 12\qquad \mathrm{(E) \ } 15$


	\item (\textit{AMC 10 10B 2004}) A standard six-sided die is rolled, and $P$ is the product of the five numbers that are visible. What is the largest number that is certain to divide $P$?


 $\mathrm{(A) \ } 6 \qquad \mathrm{(B) \ } 12 \qquad \mathrm{(C) \ } 24 \qquad \mathrm{(D) \ } 144\qquad \mathrm{(E) \ } 720$


	\item (\textit{AMC 10 10B 2004}) In the expression $c\cdot a^b-d$, the values of $a$, $b$, $c$, and $d$ are $0$, $1$, $2$, and $3$, although not necessarily in that order. What is the maximum possible value of the result?


 $\mathrm{(A) \ } 5 \qquad \mathrm{(B) \ } 6\qquad \mathrm{(C) \ } 8 \qquad \mathrm{(D) \ } 9\qquad \mathrm{(E) \ } 10$


	\item (\textit{AMC 10 10B 2004}) Which of the following numbers is a perfect square?


 $\noindent \textbf{(A)}\ 98! \cdot 99! \qquad \textbf{(B)}\ 98! \cdot 100! \qquad \textbf{(C)}\ 99! \cdot 100! \qquad \textbf{(D)}\ 99! \cdot 101! \qquad \textbf{(E)}\ 100! \cdot 101!$


	\item (\textit{AMC 10 10B 2004}) On a trip from the United States to Canada, Isabella took $d$ U.S. dollars. At the border she exchanged them all, receiving $10$ Canadian dollars for every $7$ U.S. dollars. After spending $60$ Canadian dollars, she had $d$ Canadian dollars left. What is the sum of the digits of $d$?


 $\mathrm{(A) \ } 5 \qquad \mathrm{(B) \ } 6\qquad \mathrm{(C) \ } 7\qquad \mathrm{(D) \ } 8\qquad \mathrm{(E) \ } 9$


	\item (\textit{AMC 10 10B 2004}) Minneapolis-St. Paul International Airport is $8$ miles southwest of downtown St. Paul and $10$ miles southeast of downtown Minneapolis. Which of the following is closest to the number of miles between downtown St. Paul and downtown Minneapolis?


 $\mathrm{(A) \ } 13 \qquad \mathrm{(B) \ } 14\qquad \mathrm{(C) \ } 15\qquad \mathrm{(D) \ } 16\qquad \mathrm{(E) \ } 17$


	\item (\textit{AMC 10 10B 2004}) A square has sides of length $10$, and a circle centered at one of its vertices has radius $10$. What is the area of the union of the regions enclosed by the square and the circle?


 $\mathrm{(A) \ } 200+25\pi \qquad \mathrm{(B) \ } 100+75\pi \qquad \mathrm{(C) \ } 75+100\pi \qquad \mathrm{(D) \ } 100+100\pi \qquad \mathrm{(E) \ } 100+125\pi$


	\item (\textit{AMC 10 10B 2004}) A grocer makes a display of cans in which the top row has one can and each lower row has two more cans than the row above it. If the display contains $100$ cans, how many rows does it contain?


 $\mathrm{(A) \ } 5 \qquad \mathrm{(B) \ } 8 \qquad \mathrm{(C) \ } 9 \qquad \mathrm{(D) \ } 10\qquad \mathrm{(E) \ } 11$


	\item (\textit{AMC 10 10B 2004}) Two eight-sided dice each have faces numbered $1$ through $8$. When the dice are rolled, each face has an equal probability of appearing on the top. What is the probability that the product of the two top numbers is greater than their sum?


 $\mathrm{(A) \ } \frac{1}{2} \qquad \mathrm{(B) \ } \frac{47}{64} \qquad \mathrm{(C) \ } \frac{3}{4} \qquad \mathrm{(D) \ } \frac{55}{64} \qquad \mathrm{(E) \ } \frac{7}{8}$


	\item (\textit{AMC 10 10B 2004}) An annulus is the region between two concentric circles. The concentric circles in the figure have radii $b$ and $c$, with $b>c$. Let $OX$ be a radius of the larger circle, let $XZ$ be tangent to the smaller circle at $Z$, and let $OY$ be the radius of the larger circle that contains $Z$. Let $a=XZ$, $d=YZ$, and $e=XY$. What is the area of the annulus? 


 
                    \begin{figure}[H]
                        \centering
                        \includegraphics[width=0.5\textwidth]{Problem/AMC_10_10B_2004_12_Problem_Figure_1.png}
                    \end{figure}
                    
$\mathrm{(A) \ } \pi a^2 \qquad \mathrm{(B) \ } \pi b^2 \qquad \mathrm{(C) \ } \pi c^2 \qquad \mathrm{(D) \ } \pi d^2 \qquad \mathrm{(E) \ } \pi e^2$


	\item (\textit{AMC 10 10B 2004}) In the United States, coins have the following thicknesses: penny, $1.55$ mm; nickel, $1.95$ mm; dime, $1.35$ mm; quarter, $1.75$ mm. If a stack of these coins is exactly $14$ mm high, how many coins are in the stack?


 $\mathrm{(A) \ } 7 \qquad \mathrm{(B) \ } 8 \qquad \mathrm{(C) \ } 9 \qquad \mathrm{(D) \ } 10 \qquad \mathrm{(E) \ } 11$


	\item (\textit{AMC 10 10B 2004}) A bag initially contains red marbles and blue marbles only, with more blue than red. Red marbles are added to the bag until only $\frac{1}{3}$ of the marbles in the bag are blue. Then yellow marbles are added to the bag until only $\frac{1}{5}$ of the marbles in the bag are blue. Finally, the number of blue marbles in the bag is doubled. What fraction of the marbles now in the bag are blue?


 $\mathrm{(A) \ } \frac{1}{5} \qquad \mathrm{(B) \ } \frac{1}{4} \qquad \mathrm{(C) \ } \frac{1}{3} \qquad \mathrm{(D) \ } \frac{2}{5} \qquad \mathrm{(E) \ } \frac{1}{2}$


	\item (\textit{AMC 10 10B 2004}) Patty has $20$ coins consisting of nickels and dimes. If her nickels were dimes and her dimes were nickels, she would have $70$ cents more. How much are her coins worth?


 $\textbf{(A)}\ \textdollar 1.15\qquad\textbf{(B)}\ \textdollar 1.20\qquad\textbf{(C)}\ \textdollar 1.25\qquad\textbf{(D)}\ \textdollar 1.30\qquad\textbf{(E)}\ \textdollar 1.35$


	\item (\textit{AMC 10 10B 2004}) Three circles of radius $1$ are externally tangent to each other and internally tangent to a larger circle. What is the radius of the large circle?


 $\mathrm{(A) \ } \frac{2 + \sqrt{6}}{3} \qquad \mathrm{(B) \ } 2 \qquad \mathrm{(C) \ } \frac{2 + 3\sqrt{2}}{2} \qquad \mathrm{(D) \ } \frac{3 + 2\sqrt{3}}{3} \qquad \mathrm{(E) \ } \frac{3 + \sqrt{3}}{2}$


	\item (\textit{AMC 10 10B 2004}) The two digits in Jack's age are the same as the digits in Bill's age, but in reverse order. In five years Jack will be twice as old as Bill will be then. What is the difference in their current ages?


 $\mathrm{(A) \ } 9 \qquad \mathrm{(B) \ } 18 \qquad \mathrm{(C) \ } 27 \qquad \mathrm{(D) \ } 36\qquad \mathrm{(E) \ } 45$


	\item (\textit{AMC 10 10B 2004}) In the right triangle $\triangle ACE$, we have $AC=12$, $CE=16$, and $EA=20$. Points $B$, $D$, and $F$ are located on $AC$, $CE$, and $EA$, respectively, so that $AB=3$, $CD=4$, and $EF=5$. What is the ratio of the area of $\triangle DBF$ to that of $\triangle ACE$?


 
                    \begin{figure}[H]
                        \centering
                        \includegraphics[width=0.5\textwidth]{Problem/AMC_10_10B_2004_18_Problem_Figure_1.png}
                    \end{figure}
                    
$\mathrm{(A) \ } \frac{1}{4} \qquad \mathrm{(B) \ } \frac{9}{25} \qquad \mathrm{(C) \ } \frac{3}{8} \qquad \mathrm{(D) \ } \frac{11}{25} \qquad \mathrm{(E) \ } \frac{7}{16}$


	\item (\textit{AMC 10 10B 2004}) In the sequence $2001$, $2002$, $2003$, $\ldots$ , each term after the third is found by subtracting the previous term from the sum of the two terms that precede that term. For example, the fourth term is $2001 + 2002 - 2003 = 2000$. What is the
$2004^\textrm{th}$ term in this sequence?


 $\mathrm{(A) \ } -2004 \qquad \mathrm{(B) \ } -2 \qquad \mathrm{(C) \ } 0 \qquad \mathrm{(D) \ } 4003 \qquad \mathrm{(E) \ } 6007$


	\item (\textit{AMC 10 10B 2004}) In $\triangle ABC$ points $D$ and $E$ lie on $BC$ and $AC$, respectively. If $AD$ and $BE$ intersect at $T$ so that $\frac{AT}{DT}=3$ and $\frac{BT}{ET}=4$, what is $\frac{CD}{BD}$?


 
                    \begin{figure}[H]
                        \centering
                        \includegraphics[width=0.5\textwidth]{Problem/AMC_10_10B_2004_20_Problem_Figure_1.png}
                    \end{figure}
                    
$\mathrm{(A) \ } \frac{1}{8} \qquad \mathrm{(B) \ } \frac{2}{9} \qquad \mathrm{(C) \ } \frac{3}{10} \qquad \mathrm{(D) \ } \frac{4}{11} \qquad \mathrm{(E) \ } \frac{5}{12}$


	\item (\textit{AMC 10 10B 2004}) Let $1$; $4$; $\ldots$ and $9$; $16$; $\ldots$ be two arithmetic progressions. The set $S$ is the union of the first $2004$ terms of each sequence. How many distinct numbers are in $S$?


 $\mathrm{(A) \ } 3722 \qquad \mathrm{(B) \ } 3732 \qquad \mathrm{(C) \ } 3914 \qquad \mathrm{(D) \ } 3924 \qquad \mathrm{(E) \ } 4007$


	\item (\textit{AMC 10 10B 2004}) A triangle with sides of length $5, 12,$ and $13$ has both an inscribed and a circumscribed circle. What is the distance between the centers of those circles?


 $\mathrm{(A) \ } \frac{3\sqrt{5}}{2} \qquad \mathrm{(B) \ } \frac{7}{2} \qquad \mathrm{(C) \ } \sqrt{15} \qquad \mathrm{(D) \ } \frac{\sqrt{65}}{2} \qquad \mathrm{(E) \ } \frac{9}{2}$


	\item (\textit{AMC 10 10B 2004}) Each face of a cube is painted either red or blue, each with probability $1/2$. The color of each face is determined independently. What is the probability that the painted cube can be placed on a horizontal surface so that the four vertical faces are all the same color?


 $\mathrm{(A) \ } \frac{1}{4} \qquad \mathrm{(B) \ } \frac{5}{16} \qquad \mathrm{(C) \ } \frac{3}{8} \qquad \mathrm{(D) \ } \frac{7}{16} \qquad \mathrm{(E) \ } \frac{1}{2}$


	\item (\textit{AMC 10 10B 2004}) In $\bigtriangleup ABC$ we have $AB = 7$, $AC = 8$, and $BC = 9$. Point $D$ is on the circumscribed circle of the triangle so that $AD$ bisects $\angle BAC$. What is the value of $\frac{AD}{CD}$?


 $\mathrm{(A) \ } \frac{9}{8} \qquad \mathrm{(B) \ } \frac{5}{3} \qquad \mathrm{(C) \ } 2 \qquad \mathrm{(D) \ } \frac{17}{7} \qquad \mathrm{(E) \ } \frac{5}{2}$


	\item (\textit{AMC 10 10B 2004}) A circle of radius $1$ is internally tangent to two circles of radius $2$ at points $A$ and $B$, where $AB$ is a diameter of the smaller circle. What is the area of the region, shaded in the picture, that is outside the smaller circle and inside each of the two larger circles? 


 $\mathrm{(A) \ } \frac{5}{3} \pi - 3\sqrt 2 \qquad  \mathrm{(B) \ } \frac{5}{3} \pi - 2\sqrt 3 \qquad  \mathrm{(C) \ } \frac{8}{3} \pi - 3\sqrt 3 \qquad  \mathrm{(D) \ } \frac{8}{3} \pi - 3\sqrt 2 \qquad  \mathrm{(E) \ } \frac{8}{3} \pi - 2\sqrt 3$

                    \begin{figure}[H]
                        \centering
                        \includegraphics[width=0.5\textwidth]{Problem/AMC_10_10B_2004_25_Problem_Figure_1.png}
                    \end{figure}
                    


\end{enumerate}